Three aspects of contractual implementation follow from this accounting. First, the state variable used by a trigger must be informative about repayment resources at the date service is due. A timely and verifiable index may be preferable to a broader economic measure that becomes available only after a financing crisis has developed, even if the latter is more closely related to annual output. Second, the size of relief has to be assessed jointly with the promise made outside the trigger. A high activation frequency is not evidence of good insurance when most activations occur in states that would already have repaid. Third, a suspension should be evaluated together with its capitalization and repayment schedule. Moving an amount out of the current budget does not establish that it has disappeared from the sovereign balance sheet.

These considerations suggest a practical sequence for empirical work. The starting point would be the legal trigger and service schedule, followed by a joint distribution of the index, output, and feasible fiscal adjustment. Pricing data would then help separate expected credit losses from additional liquidity or ambiguity compensation. For a deferral, the state must also record the service generated by earlier activations; otherwise, repeated relief can be mistaken for a series of independent gifts. Observed sovereign debt stocks would have to be mapped into the model's maturity-specific obligations before the calculated limits could be compared with actual fiscal headroom.

The analytical extensions retained in the appendices provide context for that agenda. Maturity changes the dates over which a sovereign must retain market access; adaptation can improve future resources at a current fiscal cost; and private-asset prices can react to the same aggregate states through cash flows and discount rates. Monetary collateral treatment operates under its own mandate and cannot be inferred from an incompletely specified consolidated planner. These channels are economically relevant even though they are not needed to establish the main contract-design result. Keeping them explicit, with their assumptions and unresolved empirical requirements, preserves the broader setting of the fiscal-limit analysis while distinguishing it from the calculations actually reported.
